\subsection{Comparador binario}
\subsubsection{Consigna}
El circuito es un comparador binario de dos números ($A$ y $B$) de dos bits cada uno.
Las salidas ($S_0$, $S_1$ y $S_2$) representan la salida del comparador, donde:

\begin{itemize}
    \item $S_0 = 1$ cuando $A > B$.
    \item $S_1 = 1$ cuando $A < B$.
    \item $S_2 = 1$ cuando $A = B$.
\end{itemize}

En caso de no darse la condición, la salida permanece en cero.

\begin{enumerate}[nosep]
    \item \textbf{Tabla de Verdad:} Diseñar la tabla de estados completa.
    \item \textbf{Funciones Lógicas:} Obtener las funciones lógicas de salidas con circuitos combinacionales.
    \item \textbf{Minimización por Karnaugh:} Minimizar utilizando mapas de Karnaugh.
    \item \textbf{Minimización Algebraica:} Minimizar utilizando los teoremas y postulados del álgebra de Boole.
    \item \textbf{Implementación Física:} Armar el circuito y verificar su funcionamiento en el MiniLab.
    \item \textbf{Simulación Virtual:} Armar el circuito y verificar su funcionamiento en el simulador ``fals-dat.com''.
\end{enumerate}
